\documentclass[11pt]{article}

\usepackage[margin=1in]{geometry}
\usepackage{amsmath,amssymb}
\usepackage{enumitem}
\usepackage{hyperref}
\usepackage{xcolor}
\usepackage{booktabs}
\usepackage{listings}

\hypersetup{colorlinks=true, linkcolor=blue, urlcolor=blue, citecolor=blue}

\lstset{
  basicstyle=\ttfamily\small,
  breaklines=true,
  columns=fullflexible,
}

\title{TOAE209/TOAH209 -- Design and Analysis of Algorithms\\[4pt]
\large Week 12: Practical Exercises}
\author{Foreign Trade University}
\date{}

\begin{document}
\maketitle

\section*{Instructions}

For each problem below there is a starter file
\texttt{exercises/starter\_code\_problemNN.py} containing function signatures with
\texttt{raise NotImplementedError} bodies, and a small \texttt{\_\_main\_\_} block with
tests. Implement the missing functions so the tests pass. A reference solution is
provided in \texttt{solutions/solution\_problemNN.py} -- try the problem yourself before
looking!

All problems use only the Python 3.10+ standard library. Shared helper functions live in
\texttt{starter\_code.py} at the week root: \texttt{num\_variables},
\texttt{all\_assignments}, \texttt{generate\_3sat}, \texttt{make\_graph},
\texttt{generate\_graph}, \texttt{neighbors}, \texttt{complement\_graph},
\texttt{is\_clique}, \texttt{generate\_subset\_sum}, and
\texttt{generate\_distance\_matrix}.

A \emph{CNF formula} is a list of clauses, each clause a list of signed integers (a
positive literal \texttt{v} means ``variable $v$ is true,'' \texttt{-v} means false); a
\emph{graph} is a pair \texttt{(n, edges)} with \texttt{edges} a set of
\texttt{frozenset(\{u, v\})}. These problems concern \emph{intractable} (NP-hard) problems,
so the reference algorithms are deliberately exponential brute-force / backtracking
searches -- correct, but only practical for small inputs.

The 12 problems are grouped into three themes:
\begin{itemize}
    \item \textbf{Verifiers \& Brute Force} (Problems 1--5): SAT verification and brute-force
    solving; verifiers and brute-force optima for independent set, vertex cover, and clique,
    including the complementarity identities $\mathrm{VC} = V \setminus \mathrm{IS}$ and
    ``max clique in $G$ = max independent set in $\overline{G}$.''
    \item \textbf{Reductions} (Problems 6, 12): the gadget reduction \textsc{3-SAT} $\to$
    \textsc{Independent Set}, and an explicit mapping reduction between \textsc{Vertex
    Cover} and \textsc{Independent Set} with the equivalence ``$x \in A \iff f(x) \in B$''
    checked on random instances.
    \item \textbf{Search} (Problems 7--11): brute-force \textsc{Subset Sum}; backtracking
    \textsc{Hamiltonian Cycle}; exact \textsc{TSP} via Held--Karp bitmask DP (verified
    against brute force); backtracking graph $k$-coloring and the chromatic number.
\end{itemize}

\newpage
\section{Verifiers \& Brute Force}

\subsection*{Problem 1 -- SAT Verifier}

Implement \texttt{verify\_sat(formula, assignment)}: given a CNF \texttt{formula} (list of
clauses, each a list of signed ints) and an \texttt{assignment} (a tuple/list of booleans
where \texttt{assignment[v-1]} is the value of variable $v$), return \texttt{True} iff the
assignment satisfies the formula -- i.e.\ every clause has at least one true literal. A
literal \texttt{v > 0} is true iff \texttt{assignment[v-1]} is \texttt{True}; a literal
\texttt{v < 0} is true iff \texttt{assignment[-v-1]} is \texttt{False}.

\textbf{Example.} \texttt{verify\_sat([[1, -2], [2, 3]], (True, False, False)) == True}.
The empty formula \texttt{[]} is satisfied by any assignment; a formula containing the
empty clause \texttt{[[]]} is satisfied by none.

\subsection*{Problem 2 -- Brute-Force SAT Solver}

Implement \texttt{brute\_force\_sat(formula)}: try all $2^n$ assignments over the $n =
\texttt{num\_variables(formula)}$ variables (use \texttt{all\_assignments} from
\texttt{starter\_code.py}) and return the first satisfying assignment found (as a tuple of
booleans), or \texttt{None} if the formula is unsatisfiable. This is the exponential
``certificate search'' that the verifier of Problem~1 short-circuits when a certificate is
already known.

\subsection*{Problem 3 -- Independent-Set Verifier and Maximum Independent Set}

Implement \texttt{is\_independent\_set(graph, subset)}: return \texttt{True} iff no two
vertices of \texttt{subset} (an iterable of vertex labels) are joined by an edge of
\texttt{graph}. Then implement \texttt{max\_independent\_set(graph)}: by checking all $2^n$
vertex subsets (largest first), return a maximum independent set as a \texttt{frozenset}.
Intended for $n \le 18$.

\subsection*{Problem 4 -- Vertex-Cover Verifier, Minimum Vertex Cover, and Complementarity}

Implement \texttt{is\_vertex\_cover(graph, subset)}: return \texttt{True} iff every edge of
\texttt{graph} has at least one endpoint in \texttt{subset}. Implement
\texttt{min\_vertex\_cover(graph)}: by brute force, return a minimum vertex cover as a
\texttt{frozenset}. Finally implement \texttt{verify\_vc\_complement(graph)}: return
\texttt{True} iff $V \setminus (\text{a maximum independent set})$ is a minimum vertex cover
-- i.e.\ the sizes satisfy $|\mathrm{VC}| = n - |\mathrm{IS}|$ \emph{and} the complement of a
maximum independent set is in fact a valid vertex cover.

\subsection*{Problem 5 -- Clique Verifier, Maximum Clique, and the Complement Identity}

Implement \texttt{is\_clique\_subset(graph, subset)}: return \texttt{True} iff every pair of
vertices in \texttt{subset} is joined by an edge. Implement
\texttt{max\_clique(graph)}: by brute force, return a maximum clique as a \texttt{frozenset}.
Finally implement \texttt{verify\_clique\_is\_complement(graph)}: return \texttt{True} iff
the size of a maximum clique in \texttt{graph} equals the size of a maximum independent set
in \texttt{complement\_graph(graph)} (using \texttt{max\_independent\_set} from Problem~3).

\newpage
\section{Reductions}

\subsection*{Problem 6 -- Reduction \textsc{3-SAT} \texorpdfstring{$\to$}{->} Independent Set}

Implement \texttt{reduce\_3sat\_to\_independent\_set(formula)} returning \texttt{(graph, k)}:
\begin{itemize}
    \item Create one vertex per (clause, literal-position) pair, so a formula with $k$
    clauses of $\le 3$ literals each yields up to $3k$ vertices.
    \item Add \emph{triangle edges}: within each clause, join all of its literal-vertices
    pairwise.
    \item Add \emph{conflict edges}: between any two vertices in \emph{different} clauses
    whose literals are complementary (one is \texttt{v}, the other \texttt{-v}).
    \item Set $k$ = number of clauses.
\end{itemize}
Then implement \texttt{verify\_reduction(formula)}: return \texttt{True} iff
\texttt{brute\_force\_sat(formula)} finds the formula satisfiable \emph{exactly when} the
constructed graph has an independent set of size $\ge k$ (compute the latter with
\texttt{max\_independent\_set} from Problem~3). This is the both-directions correctness
check of the reduction, done computationally on small formulas.

\textbf{Tip.} You may import \texttt{brute\_force\_sat} (Problem~2) and
\texttt{max\_independent\_set} / \texttt{is\_independent\_set} (Problem~3) -- the reference
solution re-implements the small pieces it needs so each file is self-contained.

\subsection*{Problem 12 -- Reduction Certificate Demo: \textsc{Vertex Cover} \texorpdfstring{$\leftrightarrow$}{<->} \textsc{Independent Set}}

Implement a poly-time mapping reduction and verify the certificate equivalence ``$x \in A
\iff f(x) \in B$'' on random small instances.

Define the decision predicates \texttt{has\_independent\_set(graph, k)} (does \texttt{graph}
have an independent set of size $\ge k$?) and \texttt{has\_vertex\_cover(graph, k)} (size
$\le k$?), each by brute force. Implement the mapping
\texttt{reduce\_is\_to\_vc(graph, k)} returning \texttt{(graph, n - k)} (leave the graph
unchanged, complement the threshold), reflecting the identity ``$G$ has an IS of size $\ge
k$ iff $G$ has a VC of size $\le n - k$.'' Finally implement
\texttt{verify\_mapping(num\_trials, n, seed)}: for \texttt{num\_trials} random graphs and
a sweep of thresholds $k$, confirm \texttt{has\_independent\_set(graph, k)} equals
\texttt{has\_vertex\_cover(*reduce\_is\_to\_vc(graph, k))} for every $(graph, k)$. Return
\texttt{True} iff the equivalence holds in every case.

\newpage
\section{Search}

\subsection*{Problem 7 -- Subset-Sum Verifier and Brute-Force Decision}

Implement \texttt{verify\_subset\_sum(weights, subset\_indices, target)}: return
\texttt{True} iff the items of \texttt{weights} at the given indices sum to exactly
\texttt{target}. Then implement \texttt{subset\_sum\_exists(weights, target)}: by checking
all $2^n$ subsets, return \texttt{True} iff some subset of \texttt{weights} sums exactly to
\texttt{target}. Also return a witness via
\texttt{find\_subset\_sum(weights, target)}, which returns a list of indices summing to
\texttt{target}, or \texttt{None} if none exists.

\textbf{Example.} \texttt{subset\_sum\_exists([3, 34, 4, 12, 5, 2], 9) == True}
($4 + 5 = 9$).

\subsection*{Problem 8 -- Hamiltonian Cycle by Backtracking}

Implement \texttt{verify\_hamiltonian\_cycle(graph, cycle)}: given a candidate
\texttt{cycle} (a list of vertices, \emph{without} repeating the start at the end), return
\texttt{True} iff it visits every vertex exactly once and consecutive vertices (including
the wrap-around from last to first) are joined by edges. Then implement
\texttt{find\_hamiltonian\_cycle(graph)}: search by backtracking for a Hamiltonian cycle,
returning it as a list of vertices (starting at vertex 0), or \texttt{None} if the graph
has none.

\textbf{Example.} The 4-cycle $0\!-\!1\!-\!2\!-\!3\!-\!0$ has a Hamiltonian cycle; a
``path'' graph $0\!-\!1\!-\!2\!-\!3$ (no edge $3\!-\!0$) does not.

\subsection*{Problem 9 -- TSP via Held--Karp Bitmask DP}

Implement \texttt{tsp\_held\_karp(dist)}: given an $n \times n$ symmetric distance matrix,
return the length of a shortest tour visiting every city exactly once and returning to the
start (city 0), using the Held--Karp dynamic program over subsets:
\begin{verbatim}
dp[S][j] = shortest path that starts at 0, visits exactly the cities in S,
           and ends at j (where 0 in S, j in S).
answer  = min over j != 0 of dp[FULL][j] + dist[j][0].
\end{verbatim}
This runs in $O(n^2 2^n)$ time. Then implement \texttt{tsp\_brute\_force(dist)} that tries
all $(n-1)!$ permutations of the non-start cities, and
\texttt{verify\_tsp(num\_trials, n, seed)} returning \texttt{True} iff Held--Karp and brute
force agree (within $10^{-9}$) on \texttt{num\_trials} random matrices.

\subsection*{Problem 10 -- Graph $k$-Coloring by Backtracking}

Implement \texttt{verify\_coloring(graph, coloring)}: given \texttt{coloring} (a list/dict
mapping each vertex to a color index), return \texttt{True} iff it uses at most the intended
colors and no edge joins two equally-colored vertices. Then implement
\texttt{k\_coloring(graph, k)}: search by backtracking for a proper coloring of
\texttt{graph} with colors $\{0, \ldots, k-1\}$, returning it as a list (indexed by vertex)
or \texttt{None} if no proper $k$-coloring exists.

\textbf{Example.} A triangle $K_3$ needs 3 colors: \texttt{k\_coloring} returns
\texttt{None} for $k = 2$ and a valid coloring for $k = 3$.

\subsection*{Problem 11 -- Chromatic Number (Trying $k = 1, 2, 3, \ldots$)}

Implement \texttt{chromatic\_number(graph)}: using \texttt{k\_coloring} from Problem~10,
return the smallest $k \ge 0$ for which \texttt{graph} has a proper $k$-coloring (a graph
with no edges and $\ge 1$ vertex needs $1$ color; the empty graph on $0$ vertices needs
$0$). Then implement \texttt{is\_three\_colorable(graph)} returning \texttt{True} iff
$\texttt{chromatic\_number(graph)} \le 3$. \emph{Sanity check (reduction intuition):} a
bipartite graph (e.g.\ an even cycle) has chromatic number $\le 2$, while an odd cycle has
chromatic number $3$.

\newpage
\section{LeetCode Practice}

After completing the problems above, practise this week's intractability techniques
(verification, reductions, and exponential backtracking / bitmask DP for NP-hard problems)
on the following \href{https://leetcode.com}{LeetCode} problems. NP-complete problems are
attacked here with exponential backtracking / branch-and-bound. Difficulty is LeetCode's
own label.

\begin{itemize}
  \item \href{https://leetcode.com/problems/subsets/}{Subsets} -- \emph{Medium} -- Backtracking.
  \item \href{https://leetcode.com/problems/combination-sum/}{Combination Sum} -- \emph{Medium} -- Backtracking.
  \item \href{https://leetcode.com/problems/palindrome-partitioning/}{Palindrome Partitioning} -- \emph{Medium} -- Backtracking.
  \item \href{https://leetcode.com/problems/word-search/}{Word Search} -- \emph{Medium} -- Backtracking.
  \item \href{https://leetcode.com/problems/n-queens/}{N-Queens} -- \emph{Hard} -- Backtracking.
  \item \href{https://leetcode.com/problems/sudoku-solver/}{Sudoku Solver} -- \emph{Hard} -- Constraint backtracking.
  \item \href{https://leetcode.com/problems/partition-to-k-equal-sum-subsets/}{Partition to K Equal Sum Subsets} -- \emph{Medium} -- NP-hard (subset-sum).
  \item \href{https://leetcode.com/problems/flower-planting-with-no-adjacent/}{Flower Planting With No Adjacent} -- \emph{Medium} -- Graph coloring.
  \item \href{https://leetcode.com/problems/find-the-shortest-superstring/}{Find the Shortest Superstring} -- \emph{Hard} -- TSP / bitmask DP.
\end{itemize}

\end{document}
